%------------------------------------------------------------------------------%
\section{Usos da Série de Taylor}

Nessa seção apresentamos algumas exemplos de uso das Séries de Taylor para
resolver problemas matemáticos.

%------------------------------------------------------------------------------%
\subsection{Avaliando Integrais Não Elementares}

Podemos utilizar as Séries de Taylor para calcular integrais de funções
cuja integração seja muito complexa ou que a primitiva não pode ser expressa
como uma combinação finita de funções elementares.
Para ilustrar essa técnica vamos aplicá-la para uma função

\begin{example}
  Calcule a Série de Taylor para a primitiva da função
  \[
    f(x) = \sin\left(x^2\right) \;.
  \]
  \solution
  Sabemos que a função seno é igual a sua Série de Taylor para todo número real,
  isso é, podemos escrever
  \[
    \sin(\alpha) = \alpha -\frac{\alpha^3}{3!} +\frac{\alpha^5}{5!} -\frac{\alpha^7}{7!} +\cdots
  \]
  Fazendo \(\alpha = x^2 \) temos
  \[
    \sin\left(x^2\right) = x^2 -\frac{x^6}{3!} +\frac{x^{10}}{5!} -\frac{x^{14}}{7!} + \cdots
  \]
  Substituindo a função pela sua série dentro da integral e integrando termo a termo, obtemos
  \begin{align*}
    \int \sin \left(x^2\right) \,dx & = \int x^2
      - \frac{x^6}{3!}
      + \frac{x^{10}}{5!}
      - \frac{x^{14}}{7!}
      + \cdots  \,dx                                    \\[5mm]
                                    & = \int x^2 \,dx
      - \int \frac{x^6}{3!} \,dx
      + \int \frac{x^{10}}{5!} \,dx
      - \int \frac{x^{14}}{7!} \,dx
      + \cdots       \\[5mm]
                                    & = C
      + \frac{x^3}{3}
      - \frac{x^7}{7\cdot 3!}
      + \frac{x^{11}}{11\cdot 5!}
      - \frac{x^{15}}{15\cdot 7!}
      + \cdots
  \end{align*}
  onde $C$ é a constante de integração.
  Para escrevermos essa série em um somatório, colocamos $x^3$
  em evidência e escrevemos as potências em termos de $x^4$
  \begin{align*}
    \int \sin \left(x^2\right) \,dx
    & = x^3 \left(
        \frac{1}{3}
      - \frac{x^4}{7\cdot 3!}
      + \frac{x^8}{11\cdot 5!}
      - \frac{x^{12}}{15\cdot 7!}
      + \cdots
      \right) \\[3mm]
    & = x^3 \left(
        \frac{\left(-x^4\right)^0}{3}
      + \frac{(-x^4)^1}{7\cdot 3!}
      + \frac{\left(-x^4\right)^2}{11\cdot 5!}
      + \frac{\left(-x^4\right)^3}{15\cdot 7!}
      + \cdots
      \right)
  \end{align*}
  onde escolhemos $C=0$ para simplificar as expressões.
  Precisamos agora identificar as sequências que constroem os valores necessários
  \[
    \begin{array}{w{c}{2cm}w{c}{3cm}w{c}{3cm}}
      \hline
      k & 3+4k & 1+2k \\ \hline
      0 & 3    & 1    \\
      1 & 7    & 3    \\
      2 & 11   & 5    \\
      3 & 15   & 7    \\ \hline
    \end{array}
  \]
  Temos então que
  \[
    \int \sin \left(x^2\right) \,dx
    = x^3 \sum_{k=0}^\infty  \frac{\left(-x^4\right)^k}{\, (3+4k) (1+2k)! \,} \;.
  \]
\end{example}

Note que continuamos sem uma expressão fechada e finita para a integral. Porém,
mesmo a série sendo uma representação infinita, ela pode ser utilizada para aproximar
o valor da integral em qualquer ponto com a precisão que desejarmos.

%------------------------------------------------------------------------------%
\subsection{Avaliando Formas Indeterminadas}

Outra operação que pode ser complexa é o cálculo de limites indeterminados,
nesse caso a Série de Taylor também pode ajudar.
Vamos ilustrar a técnica caculando o limite
\[
  \lim_{x\to 1} \frac{\ln x}{\;x-1\;} \;.
\]
Esse é um limite, com indeterminação do tipo $\sfrac{0}{0}$,
onde podemos utilizar a regra de l'Hôpital
\[
  \lim_{x\to 1} \frac{\ln x}{\;x-1\;}
  = \lim_{x\to 1} \frac{\;\sfrac{1}{x}\;}{1}
  = 1 \;.
\]
Entretanto, a Série de Taylor oferece uma forma alternativa
para calcular esse limite que pode ser útil em outras situações.

\begin{example}
  Calcule o limite
  \(
    \lim_{x\to 1} \frac{\ln x}{\;x-1\;}
  \).
  \solution
  Primeiro obtemos a Série de Taylor da função $\ln(x)$ com centro em $a=1$
  \[
    \ln(x) = (x-1) - \frac{(x-1)^2}{2}
    + \frac{(x-1)^3}{3}
    - \frac{(x-1)^4}{4}
    + \cdots
  \]
  a igualdade é verdadeira para todo $x$ tal que $\abs{x-1} < 1$.
  Agora substituímos o logaritmo pela sua série na função original
  \[
    \frac{\ln x}{\,x-1\,}
    = \frac{1}{\,x-1\,}
      \left( (x-1)
        - \frac{(x-1)^2}{2}
        + \frac{(x-1)^3}{3}
        - \frac{(x-1)^4}{4}
        + \cdots
      \right) \;.
  \]
  Manipulamos os termos da série obtemos
  \[
    \frac{\ln x}{\,x-1\,}
    = 1 - \frac{x-1}{2} + \frac{(x-1)^2}{3} - \frac{(x-1)^3}{4} + \cdots
  \]
  Agora calculamos o limite desejado
  \[
    \lim_{x\to 1} \frac{\ln x}{\,x-1\,}
    = \lim_{x\to 1} \left( 1
                      - \frac{x-1}{2}
                      + \frac{(x-1)^2}{3}
                      - \frac{(x-1)^3}{4}
                      + \cdots
                    \right)
  \]
  como todos os termos da série, com exceção do primeiro, vão para zero,
  quando $x\to 1$, temos que
  \[
    \lim_{x\to 1} \frac{\ln x}{\,x-1\,} = 1 \;.
  \]
\end{example}

%------------------------------------------------------------------------------%
\subsection{Definição de Novas Funções}

Podemos definir novas funções como a soma de uma Série de Taylor, essa é uma
forma para definir as funções reais conhecidas para variáveis complexas,
como ilustrado na próxima subseção,
e também é como definimos várias funções especiais. Como exemplo vamos
definir uma das \emph{Funções de Bessel}\index{Função!de Bessel},
que são as soluções da equação de Bessel
\begin{equation}
  \label{eq:bessel}
  x^2 y'' + xy' + (x^2-\alpha^2)y = 0 \;.
\end{equation}

%------------------------------------------------------------------------------%
\begin{example}
  Exiba a Função de Bessel de Primeira Espécie de Ordem Zero, $J_0(x)$.
  \solution
  Essa função é uma solução da equação de Bessel~\eqref{eq:bessel} com
  $\alpha=0$
  \[
    \mathcal{L}[y] = x^2 y'' + xy' + x^2y = 0 \;.
  \]
  Vamos buscar a solução com a forma
  \[
    y(x) = \sum_{n=0}^\infty a_n x^n \;.
  \]
  Calculando as derivadas de $y$ obtemos
  \begin{align*}
    y'(x)  & = \sum_{n=1}^\infty a_n n x^{n-1} \\[3mm]
    y''(x) & = \sum_{n=2}^\infty a_n n(n-1) x^{n-2} \;.
  \end{align*}
  Substituindo na equação temos
  \begin{align*}
    \mathcal{L}[y] & = x^2 y'' + xy' + x^2y \\[3mm]
                   & = x^2 \sum_{n=2}^\infty a_n n(n-1) x^{n-2}
                     + x   \sum_{n=1}^\infty a_n n x^{n-1}
                     + x^2 \sum_{n=0}^\infty a_n x^n
                   \\[3mm]
                   & = \sum_{n=2}^\infty a_n n (n-1) x^n
                     + \sum_{n=1}^\infty a_n n       x^n
                     + \sum_{n=0}^\infty a_n         x^{n+2} \;.
  \end{align*}
  Queremos agrupar todos os termos em um único somatório,
  para isso fazemos uma mudança de índice no último
  para igualar as potências de $x$
  \[
  \mathcal{L}[y] = \sum_{n=2}^\infty a_n n (n-1) x^n
                 + \sum_{n=1}^\infty a_n n       x^n
                 + \sum_{n=2}^\infty a_{n-2}     x^n \;.
  \]
  Precisamos agora fazer o índice do segundo somatório começar em $n=2$,
  conseguimos isso movendo o termo $n=1$ para fora
  \begin{align*}
    \mathcal{L}[y] & =        \sum_{n=2}^\infty a_n n (n-1) x^n
                     + a_1x + \sum_{n=2}^\infty a_n n       x^n
                     +        \sum_{n=2}^\infty a_{n-2}     x^n \\[3mm]
                   & = a_1x + \sum_{n=2}^\infty \big( a_n n (n-1) + a_n n + a_{n-2} \big) x^n \\[3mm]
                   & = a_1x + \sum_{n=2}^\infty \big( a_n n^2 + a_{n-2} \big) x^n \;.
  \end{align*}
  Para que
  \(
    \mathcal{L}[y] = 0
  \)
  precisamos que o coeficiente
  de cada potência de $x$ seja zero. Ao analisarmos o
  coeficiente de $x$ verificamos que
  \[
    a_1 = 0 \;.
  \]
  Para as demais potências, $x^n$, temos a condição
  \[
    a_n n^2 + a_{n-2} = 0
  \]
  que produz a relação de recorrência
  \[
    a_n = -\frac{a_{n-2}}{n^2} \;.
  \]
  Usando essa relação e a condição $a_1 = 0$ concluímos que todos
  os coeficientes ímpares são nulos
  \[
    a_n = 0 \qquad n = 1, 3, 5, \dots
  \]
  Para encontrarmos os coeficientes pares fazemos $n=2k$
  \[
    a_{2k} = -\frac{a_{2k-2}}{(2k)^2} \qquad k = 1, 2, 3, \dots
  \]
  Observando que
  \begin{align*}
    a_2 & = -\frac{a_0}{2^2} \\[3mm]
    a_4 & = -\frac{a_2}{4^2}
          = -\frac{1}{4^2}\left(-\frac{a_0}{2^2}\right)
          =  \frac{a_0}{2^6} \\[3mm]
    a_6 & = -\frac{a_4}{6^2}
          = -\frac{1}{6^2}\left(\frac{a_0}{2^6}\right)
          = -\frac{a_0}{2^6(3\cdot 2)^2} \\[3mm]
    a_8 & = -\frac{a_6}{8^2}
          = -\frac{1}{8^2}\left(-\frac{a_0}{2^6(3\cdot 2)^2}\right)
          =  \frac{a_0}{2^8(4\cdot 3\cdot 2)^2} \;.
  \end{align*}
  Generalizando temos
  \[
    a_{2k} = \frac{(-1)^k}{\,2^{2k}\,(k!)^2\,} a_0 \qquad k = 1, 2, 3, \dots
  \]
  Podemos agora escrever
  \[
    y(x) = a_0 \sum_{k=0}^\infty \frac{\,(-1)^k\,}{\,2^{2k}\,(k!)^2\,} \, x^{2k} \;.
  \]
  A \emph{Função de Bessel de Primeira Espécie de Ordem Zero} é definida como
  \[
    J_0(x) = \sum_{k=0}^\infty \frac{\,(-1)^k\,}{\,2^{2k}\,(k!)^2\,} \, x^{2k} \;.
  \]
  %----------------------------------------------------------------------------%
  Vamos agora analisar a convergência dessa série usando o Teste da Razão
  \[
    \abs{\frac{u_{k+1}}{u_k}}
     = \abs{ \frac{\,(-1)^{k+1} x^{2k+2}\,}{\,2^{2k+2}\,((k+1)!)^2\,}\;
              \frac{\,2^{2k}\,(k!)^2\,}{\,(-1)^k x^{2k}\,} }
     = \frac{x^2}{\,2^2 (k+1)^2\,}
  \]
  Quando $k\to\infty$ essa expressão vai para zero independentemente do valor
  de $x$, portanto a série converge para qualquer $x\in\R$.
  %----------------------------------------------------------------------------%
\end{example}
%------------------------------------------------------------------------------%

\begin{figure}
  \begin{center}%
    \vspace{-0.2\baselineskip}%
    \scalebox{0.9}{% Title: gl2ps_renderer figure
% Creator: GL2PS 1.4.2, (C) 1999-2020 C. Geuzaine
% For: Octave
% CreationDate: Tue Sep 28 14:13:47 2021
\setlength{\unitlength}{1pt}
\begin{picture}(0,0)
\includegraphics[width=17cm,height=7cm]{fig_taylor_bessel}
\end{picture}%
\begin{picture}(481,198)(0,0)
\fontsize{10}{0}\selectfont\put(88.9306,67.8458){\makebox(0,0)[t]{\textcolor[rgb]{0.15,0.15,0.15}{{5}}}}
\fontsize{10}{0}\selectfont\put(159.024,67.8458){\makebox(0,0)[t]{\textcolor[rgb]{0.15,0.15,0.15}{{10}}}}
\fontsize{10}{0}\selectfont\put(229.117,67.8458){\makebox(0,0)[t]{\textcolor[rgb]{0.15,0.15,0.15}{{15}}}}
\fontsize{10}{0}\selectfont\put(299.21,67.8458){\makebox(0,0)[t]{\textcolor[rgb]{0.15,0.15,0.15}{{20}}}}
\fontsize{10}{0}\selectfont\put(369.303,67.8458){\makebox(0,0)[t]{\textcolor[rgb]{0.15,0.15,0.15}{{25}}}}
\fontsize{10}{0}\selectfont\put(439.396,67.8458){\makebox(0,0)[t]{\textcolor[rgb]{0.15,0.15,0.15}{{30}}}}
\fontsize{10}{0}\selectfont\put(13.8417,20.4678){\makebox(0,0)[r]{\textcolor[rgb]{0.15,0.15,0.15}{{-0,5}}}}
\fontsize{10}{0}\selectfont\put(13.8417,130.185){\makebox(0,0)[r]{\textcolor[rgb]{0.15,0.15,0.15}{{0,5}}}}
\fontsize{10}{0}\selectfont\put(13.8417,185.044){\makebox(0,0)[r]{\textcolor[rgb]{0.15,0.15,0.15}{{1}}}}
\fontsize{10}{0}\selectfont\put(25.8468,190.53){\makebox(0,0)[l]{\textcolor[rgb]{0,0,0}{{$J_0(x)$}}}}
\fontsize{10}{0}\selectfont\put(166.033,163.101){\makebox(0,0)[l]{\textcolor[rgb]{0,0,0}{{$P_{20}$}}}}
\fontsize{10}{0}\selectfont\put(271.173,163.101){\makebox(0,0)[l]{\textcolor[rgb]{0,0,0}{{$P_{40}$}}}}
\fontsize{10}{0}\selectfont\put(376.312,163.101){\makebox(0,0)[l]{\textcolor[rgb]{0,0,0}{{$P_{60}$}}}}
\fontsize{10}{0}\selectfont\put(45.4729,31.4396){\makebox(0,0)[r]{\textcolor[rgb]{0,0,0}{{$P_{2}$}}}}
\fontsize{10}{0}\selectfont\put(104.351,31.4396){\makebox(0,0)[l]{\textcolor[rgb]{0,0,0}{{$P_{10}$}}}}
\fontsize{10}{0}\selectfont\put(215.098,31.4396){\makebox(0,0)[l]{\textcolor[rgb]{0,0,0}{{$P_{30}$}}}}
\fontsize{10}{0}\selectfont\put(320.238,31.4396){\makebox(0,0)[l]{\textcolor[rgb]{0,0,0}{{$P_{50}$}}}}
\fontsize{10}{0}\selectfont\put(425.377,31.4396){\makebox(0,0)[l]{\textcolor[rgb]{0,0,0}{{$P_{70}$}}}}
\end{picture}
}%
    \vspace{-\baselineskip}%
  \end{center}
  \caption{Comparação da função de Bessel com seus polinômios de Taylor.}
  \label{fig:funcao_tayplor_bessel}
\end{figure}

A Figura~\ref{fig:funcao_tayplor_bessel} a seguir
mostra o gráfico da função $J_0(x)$ e alguns dos seus Polinômios de Taylor.

%------------------------------------------------------------------------------%
\subsection{Identidade de Euler}

Com os devidos cuidados, os resultados sobre sequências e séries
valem também para os números ou funções complexos. Vamos apresentar aqui a
\emph{Identidade de Euler}\index{Identidade de Euler}\index{Euler!identidade}
\[
  e^{i\theta} = \cos\theta + i\sin\theta   \;.
\]
Na Seção~\ref{sec:NumerosComplexos} apresentamos um resumo sobre os
números complexos, aqui vamos utilizar que um número complexo
pode ser escrito como
\[
  z = a + bi
\]
onde $a$ e $b$ são números reais e $i$ é a unidade imaginária
\[
  i = \sqrt{-1}   \;.
\]
As operações de soma e produto com números complexos obedecem as mesmas
propriedades conhecidas para os números reais. Aplicando essas propriedades
podemos calcular as potências da unidade imaginária, por exemplo,
\[
  i^2 = -1, \qquad i^3 = -i, \qquad i^4 = 1, \qquad i^5 = i   \;.
\]
Vamos estender a definição das funções seno, cosseno e exponencial
para variáveis complexas, isso é, queremos avaliar essas funções para qualquer
valor complexo $z\in\C$. Fazemos isso substituindo o $x$ real pelo $z$ complexo
nas Séries de Taylor dessas funções
\begin{gather*}
  \sin(z) = z - \frac{z^3}{3!} + \frac{z^5}{5!} + \frac{z^7}{7!} - \cdots \\[3mm]
  \cos(z) = 1 - \frac{z^2}{2!} + \frac{z^4}{4!} + \frac{z^8}{8!} - \cdots \\[3mm]
  e^z = 1 + z + \frac{z^2}{2!} + \frac{z^3}{3!} + \frac{z^4}{4!} +  \cdots
\end{gather*}
Note que, estamos apresentando esse resultado sem a devida justificativa.
Para sermos rigorosos precisamos verificar que todas as operações estão bem
definidas e que as séries convergem, porém, essa analise foge ao escopo desse curso.

Para verificarmos a Identidade de Euler escrevemos Série de Taylor
da exponencial de $z=i\theta$, com $\theta\in\R$
\[
  e^{i\theta}
    = 1 + (i\theta)
    + \frac{(i\theta)^2}{2!}
    + \frac{(i\theta)^3}{3!}
    + \frac{(i\theta)^4}{4!}
    + \frac{(i\theta)^5}{5!}
    + \frac{(i\theta)^6}{6!}
    + \cdots
\]
Utilizando a propriedade de potências $(xy)^r = x^r y^r$ reescrevemos essa expressão como
\[
  e^{i\theta}
    = 1 + i\theta
    + \frac{i^2\theta^2}{2!}
    + \frac{i^3\theta^3}{3!}
    + \frac{i^4\theta^4}{4!}
    + \frac{i^5\theta^5}{5!}
    + \frac{i^6\theta^6}{6!}
    + \cdots
\]
Utilizando o cálculo das potencias de $i$ obtemos
\[
  e^{i\theta}
    = 1
    + i\theta
    -  \frac{\theta^2}{2!}
    - i\frac{\theta^3}{3!}
    +  \frac{\theta^4}{4!}
    + i\frac{\theta^5}{5!}
    -  \frac{\theta^6}{6!}
    +  \cdots
\]
Nessa expressão, os termos com potências impares contém $i$ e os termos
com potências pares não. Vamos agrupar esses dois tipos de termos
e colocar $i$ em evidência
\[
  e^{i\theta}
    = \left(
        1
        - \frac{\theta^2}{2!}
        + \frac{\theta^4}{4!}
        - \frac{\theta^6}{6!}
        + \cdots
      \right)
    + i \left(
          \theta
          - \frac{\theta^3}{3!}
          + \frac{\theta^5}{5!}
          - \frac{\theta^7}{7!}
          + \cdots
        \right)   \;.
\]
Os termos dentro do primeiro par de parênteses
são os termos da Série de Taylor do cosseno e os termos
dentro do segundo par de parênteses são a série do seno.
Podemos, então, escrever a Identidade de Euler
\[
  e^{i\theta} = \cos\theta + i\sin\theta   \;.
\]
Uma consequência interessante dessa identidade é que podemos escrever
uma relação envolvendo ``Todas as Constantes da Matemática''.
Para isso basta escolher $\theta=\pi$
\[
  e^{i\pi} \;=\; \cos\pi + i\sin\pi \;=\; -1 + i0 \;=\; -1
\]
ou seja
\[
  e^{i\pi} + 1 = 0   \;.
\]

%------------------------------------------------------------------------------%
\subsection{Cálculo da Série de Taylor para o Arco Tangente}

Vamos ilustrar uma forma alternativa para calcular a Série de Taylor da
função $\arctan(x)$ para ilustrar algumas manipulações possíveis
quando precisamos obter uma Série de Taylor.
Vamos começar, como em outros casos, calculando as derivadas da função
\begin{align*}
  \frac{d}{dx}     & \arctan(x) = \frac{1}{1+x^2}                     \\[3mm]
  \frac{d^2}{dx^2} & \arctan(x) = \frac{-2x}{(1+x^2)^2}               \\[3mm]
  \frac{d^3}{dx^3} & \arctan(x) = \frac{6x^2-2}{(1+x^2)^3}            \\[3mm]
  \frac{d^4}{dx^4} & \arctan(x) = \frac{24x(1-x^2)}{(1+x^2)^4}   \;.
\end{align*}
Não precisamos ir mais longe para perceber que esse não é um bom método,
as expressões estão se tornando cada vez mais complexas,
portanto, precisamos tentar outra abordagem.

Inicialmente, nesse desenvolvimento, não vamos nos preocupar com o
rigor matemático para obter uma candidata a Série de Taylor,
depois justificamos o processo utilizado demonstrando a
convergência da série.
Começamos observando que a derivada da função arco tangente
é igual a soma de uma Série Geométrica com $\alpha=1$ e $r=-x^2$
\[
  \frac{d}{dx} \arctan(x)
  = \frac{1}{1+x^2}
  = \frac{\alpha}{1-r}
  = 1 - x^2 + x^4 - x^6 + x^8 - \cdots
\]
Escrevemos então que
\[
  \frac{d}{dx} \arctan(x) =  1 - x^2 + x^4 - x^6 + x^8 - \cdots
\]
integrando os dois lados temos a candidata a Série de Taylor para $\arctan(x)$
\[
  \arctan(x) = x -\frac{x^3}{3} + \frac{x^5}{5} - \frac{x^7}{7} + \frac{x^9}{9} - \cdots
\]
Temos uma Série de Taylor, porém, nossos cálculos não foram cuidadosos.
Precisamos demonstrar que a série obtida converge para a função.
Faremos isso voltando para a fórmula da soma da Série Geométrica
\begin{equation}
  \label{eq:arctan_serie_geometrica}
  \frac{1}{1+x^2}
    = 1 - x^2 + x^4 - x^6 + \cdots
    + (-1)^n x^{2n}
    + (-1)^{n+1} x^{2(n+1)} + \cdots
\end{equation}
Note que para qualquer $n$ o resto dessa série, $R_n$,
é uma Série Geométrica com $r=-x^2$ e $\alpha=(-1)^{n+1} x^{2(n+1)}$
\[
  R_n =
  (-1)^{n+1} x^{2(n+1)} +
  (-1)^{n+2} x^{2(n+2)} +
  (-1)^{n+3} x^{2(n+3)} + \cdots
\]
Utilizando a fórmula da soma dos termos dessa série obtemos
\[
  R_n = \frac{(-1)^{n+1} x^{2(n+1)}}{1+x^2}    \;.
\]
Podemos então escrever a série em \eqref{eq:arctan_serie_geometrica}
como uma soma finita
\[
  \frac{1}{1+x^2}
    = 1 - x^2 + x^4 - x^6 + \cdots
    + (-1)^n x^{2n}
    + \frac{(-1)^{n+1} x^{2(n+1)}}{1+x^2}
\]
Temos agora uma expressão finita para a derivada da função arco tangente
\[
  \frac{d}{dx} \arctan(x)
    = 1 - x^2 + x^4 - x^6
    + \cdots
    + (-1)^n x^{2n}
    + \frac{(-1)^{n+1} x^{2(n+1)}}{1+x^2}
\]
Por ser uma soma finita podemos integrar termo a termo
sem preocupações com a convergência
\[
  \arctan(x) =
    x - \frac{x^3}{3}
    + \frac{x^5}{5}
    - \frac{x^7}{7}
    + \cdots
    + \frac{(-1)^n x^{2n+1}}{2n+1}
    + \int_0^x \frac{(-1)^{n+1} u^{2(n+1)}}{1+u^2} \,du   \;.
\]
Estamos agora em uma situação similar ao Teorema de Taylor~\ref{teo:teorema-taylor},
onde podemos concluir que a série converge para a função se o erro
\[
  R_n(x) = \int_0^x \frac{(-1)^{n+1} u^{2(n+1)}}{1+u^2} \,du
\]
convergir para zero, quando $n\to\infty$.

Para garantirmos que $R_n$ tende para zero, basta verificar que seu módulo tente a zero,
ou seja, queremos calcular o limite de
\begin{align*}
    \abs{R_n(x)}
    & = \abs{\int_0^x \frac{(-1)^{n+1} u^{2(n+1)}}{1+u^2} \,du} \\[3mm]
    & = \int_0^{\abs{x}} \abs{\frac{(-1)^{n+1} u^{2(n+1)}}{1+u^2}} \,du \\[3mm]
    & = \int_0^{\abs{x}} \frac{u^{2(n+1)}}{1+u^2} \,du
\end{align*}
quando $n\to\infty$. Como $1+u^2\geq 1$ temos que
\[
  \abs{R_n(x)}
    \leq \int_0^{\abs{x}} u^{2(n+1)} \,du
    = \frac{u^{2n+3}}{2n+3} \evalat{0}{\abs{x}}
    = \frac{\abs{x}^{2n+3}}{2n+3}    \;.
\]
Portanto, para todo $\abs{x} \leq 1$ sabemos que
\[
  \lim_{n\to\infty} R_n(x) = 0
\]
e podemos concluir que, para todo $x\in[-1,\,1\,]$ a Série de Taylor da função
arco tangente converge para a função
\[
  \arctan(x) = x - \frac{x^3}{3}
  + \frac{x^5}{5}
  - \frac{x^7}{7}
  + \frac{x^9}{9}
  - \cdots
\]

%------------------------------------------------------------------------------%
\subsection{Problema de Basileia}
\label{subsec:problema-Basileia}

O \emph{Problema de Basileia}\index{Problema!Basileia} consiste em
encontrar a soma da Série de Inversos do Quadrado,
apresentada na Definição~\ref{def:serie-inverso-quadrado},
\[
  \sum_{n=1}^\infty \frac{1}{\,n^2\,} = \frac{\,\pi^2}{6}    \;.
\]
A solução desse problema obtida por Euler em 1735 o tornou conhecido na comunidade
Matemática da época e é de grande importância
por estar relacionada ao estuda da Função Zeta de Riemann, $\zeta(z)$,
pois
\[
  \zeta(2) = \sum_{n=1}^\infty \frac{1}{\,n^2\,}    \;.
\]
Entre as aplicações desse resultado está o cálculo da probabilidade de dois
números escolhidos ao acaso serem primos entre si.
A verificação dos cálculos realizados por Euler usa a decomposição do seno
em um produto infinito que depende no Teorema de Fatoração de Weierstrass
para funções holomorfas
em todo o plano complexo, evidentemente além do escopo dessa disciplina.
Dessa forma, vamos apresentar apenas os passos principais desse desenvolvimento
e em seguida apresentamos uma demonstração
publicada por T. M. Apostol~\cite{Apostol-Proof-Euler-Missed} em 1983
que envolve operações mais simples.

O desenvolvimento dos cálculos de Euler começa com a Série de Taylor da função seno
\[
  \sin(x) = x
          - \frac{x^3}{3!}
          + \frac{x^5}{5!}
          - \frac{x^7}{7!}
          + \cdots
\]
dividindo ambos os lados por $x$ temos
\[
  \frac{\sin(x)}{x} = 1
                    - \frac{x^2}{3!}
                    + \frac{x^4}{5!}
                    - \frac{x^6}{7!}
                    + \cdots
\]
Usando o Teorema de Fatoração de Weierstrass nessa função, podemos escrever
\[
  \frac{\sin(x)}{x} = \prod_{n=1}^{\infty} \left(1-\frac{x^2}{\pi^2n^2}\right)
                    = \left(1-\frac{x^2}{ \pi^2}\right)
                      \left(1-\frac{x^2}{4\pi^2}\right)
                      \left(1-\frac{x^2}{9\pi^2}\right)
                      \cdots
\]
Igualando a Série de Taylor da função com sua fatoração temos
\begin{equation}
  \label{eq:fatoracao_seno}
  1
  - \frac{x^2}{3!}
  + \frac{x^4}{5!}
  - \frac{x^6}{7!}
  + \cdots
  = \left(1-\frac{x^2}{ \pi^2}\right)
    \left(1-\frac{x^2}{4\pi^2}\right)
    \left(1-\frac{x^2}{9\pi^2}\right)
    \cdots
\end{equation}
Como essa igualdade é válida sabemos que o coeficiente de cada potência
de $x$ é igual nas duas expressões.
Extraímos diretamente o coeficiente de $x^2$ da Série de Taylor
obtendo
\[
  -\frac{1}{3!} = -\frac{1}{6}    \;.
\]
Multiplicando os termos do produto e coletando apenas os
coeficientes de $x^2$ temos
\[
  -\left( \frac{1}{ \pi^2}
        + \frac{1}{4\pi^2}
        + \frac{1}{9\pi^2}
        + \cdots
  \right)
  = - \frac{1}{\,\pi^2\,} \sum_{n=1}^\infty \frac{1}{n^2}    \;.
\]
Igualando os coeficientes temos que
\[
  - \frac{1}{6}
  = - \frac{1}{\,\pi^2\,} \sum_{n=1}^\infty \frac{1}{n^2}    \;.
\]
Isolando a série chegamos ao resultado desejado
\[
  \zeta(2) = \sum_{n=1}^\infty \frac{1}{n^2} = \frac{\pi^2}{6}    \;.
\]

%------------------------------------------------------------------------------%

Apresentamos agora a \emph{demostração alternativa} proposta por Apostol,
que se baseia na avaliação da integral
\[
  I = \int_0^1\!\!\int_0^1 \frac{1}{\,1-xy\,}\, dxdy    \;.
\]
O primeiro passo é verificar que essa integral é igual a $\zeta(2)$.
Note que a fração é a soma da Série Geométrica com $\alpha=1$ e $r=xy$,
como os valores de $x$ e $y$ na integral estão entre zero e um
podemos escrever
\[
  I = \int_0^1\!\!\int_0^1 \frac{1}{\,1-xy\,} \,dxdy
    \;=\; \int_0^1\!\!\int_0^1 \sum_{n=0}^\infty x^n y^n \,dxdy    \;.
\]
Como essa série é absolutamente convergente podemos
trocar a ordem da soma com a integração
\begin{align*}
  I & = \sum_{n=0}^\infty \;\int_0^1\!\!\int_0^1 x^n y^n \,dxdy \\[3mm]
    & = \sum_{n=0}^\infty \;\int_0^1 \frac{y^n}{n+1} \,dy       \\[3mm]
    & = \sum_{n=0}^\infty \frac{1}{(n+1)^2}
      = \zeta(2)    \;.
\end{align*}
A segunda parte consiste em calcular a integral e mostrar que
\[
  I = \frac{\pi^2}{6}    \;.
\]
Aplicamos uma mudança de variáveis rotacionando os eixos
\ang{45} no sentido horário
\[
  x = \frac{\,u-v\,}{\sqrt{2}} \qquad\qquad
  y = \frac{\,u+v\,}{\sqrt{2}}
\]
de modo que
\[
  \frac{1}{1-xy} = \frac{2}{\,2-u^2+v^2\,}
\]
e a nova região de integração seja o quadrado, $\Omega$,
com vértices opostos em $(0, 0)$ e $\left(\sqrt{2}, 0\right)$.
O Jacobiano dessa transformação é
\[
  J = \frac{\,\partial x\,}{\partial u} \frac{\,\partial y\,}{\partial v}
    - \frac{\,\partial x\,}{\partial v} \frac{\,\partial y\,}{\partial u}
    = \frac{ 1}{\,\sqrt{2}\,} \frac{1}{\,\sqrt{2}\,}
    - \frac{-1}{\,\sqrt{2}\,} \frac{1}{\,\sqrt{2}\,}
    = 1    \;.
\]
Assim a integral assume a forma
\[
  I = \iint_\Omega \frac{2}{\,2-u^2+v^2\,} \,dudv    \;.
\]
Usando a simetria de reflexão do quadrado em relação
ao eixo $u$ podemos escrever a integral como
\[
  I = 4 \int_{0}^{\frac{1}{\sqrt{2}}}
        \int_{0}^{u}
        \frac{1}{\,2-u^2+v^2\,} \,dvdu
    + 4 \int_{\frac{1}{\sqrt{2}}}^{\sqrt{2}}
        \int_{0}^{\sqrt{2}-u}
        \frac{1}{\,2-u^2+v^2\,} \,dvdu    \;.
\]
Usando
\[
  \int_0^x \frac{dt}{a^2+t^2} = \frac{1}{\,a\,} \arctan\left(\frac{\,x\,}{a}\right)
\]
temos que
\begin{align*}
  \int_{0}^{u} \frac{dv}{\,2-u^2+v^2\,}
  & = \frac{1}{\,\sqrt{2-u^2}\,} \arctan\left(\frac{u}{\,\sqrt{2-u^2}\,}\right) \;, \\[3mm]
  \int_{0}^{\sqrt{2}-u} \frac{dv}{\,2-u^2+v^2\,}
  & = \frac{1}{\,\sqrt{2-u^2}\,} \arctan\left(\frac{\sqrt{2}-u}{\,\sqrt{2-u^2}\,}\right) \;.
\end{align*}
Voltando para o cálculo de $I$ fazemos $I=I_1+I_2$ onde
\begin{align*}
  I_1 & = 4 \int_{0}^{\frac{1}{\sqrt{2}}}
            \arctan\left(\frac{u}{\,\sqrt{2-u^2}\,}\right)
            \frac{du}{\,\sqrt{2-u^2}\,}
            \\[3mm]
  I_2 & = 4 \int_{\frac{1}{\sqrt{2}}}^{\sqrt{2}}
            \arctan\left(\frac{\sqrt{2}-u}{\,\sqrt{2-u^2}\,}\right)
            \frac{du}{\,\sqrt{2-u^2}\,}
\end{align*}
Para calcularmos $I_1$ usamos a transformação
\[
  u = \sqrt{2} \sin\left(\theta\right)
  \qquad\qquad
  du = \sqrt{2} \cos\left(\theta\right) \,d\theta
     = \sqrt{2-u^2}\,d\theta
\]
que produz
\begin{align*}
  I_1 & = 4 \int_{0}^{\frac{\pi}{6}}
            \arctan\left(\dfrac{  \sqrt{2} \sin\left(\theta\right)}
                               {\,\sqrt{2} \cos\left(\theta\right)\,}\right)
            \,d\theta
            \\[3mm]
      & = 4 \int_{0}^{\frac{\pi}{6}} \theta \,d\theta
        = 2 \left(\frac{\pi}{6}\right)^2
        = \frac{\,2 \pi^2}{6^2}    \;.
\end{align*}
Para $I_2$ usamos \(u = \sqrt{2} \cos\left(2\theta\right)\) que produz
\begin{align*}
  du & = -2\sqrt{2} \sin\left(2\theta\right)                \,d\theta \\[1mm]
     & = -2\sqrt{2} \sqrt{1 - \cos^2\left(2\theta\right)\,} \,d\theta \\[1mm]
     & = -2\sqrt{2} \sqrt{1 - \sfrac{u^2}{2}\,}             \,d\theta \\[1mm]
     & = -2\sqrt{2 - u^2\,}                                 \,d\theta
\end{align*}
e também
\begin{align*}
  \frac{\sqrt{2}-u}{\sqrt{2-u^2\,}}
  & = \frac{\sqrt{2}(1-\cos(2\theta))}{\sqrt{2 - 2\cos^2(2\theta)}} \\[3mm]
  & = \sqrt{\frac{1-\cos(2\theta)}{1 + \cos(2\theta)}\,}            \\[3mm]
  & = \sqrt{\frac{2\sin^2(\theta)}{\cos^2(\theta)}\,}               \\[3mm]
  & = \tan(\theta)  \;.
\end{align*}
Temos então
\[
  I_2   =   4 \int_{\frac{\pi}{6}}^{0} \arctan\left(\tan(\theta)\right) (-2)\,d\theta
      \;=\; 8 \int_{0}^{\frac{\pi}{6}} \theta \,d\theta
      \;=\; 4 \left(\frac{\pi}{6}\right)^2
      \;=\; \frac{\,4 \pi^2}{6^2}  \;.
\]
Calculando $I$ temos
\[
  I = I_1 + I_2
    = \frac{\,2 \pi^2}{6^2}
    + \frac{\,4 \pi^2}{6^2}
    = \frac{\pi^2}{6}
\]
o que conclui nossa demostração.

%------------------------------------------------------------------------------%
